\section{Basic Definitions and Examples}

\begin{definition}
    Let $R$ be a ring with $1$, and let $M$ be an abelian group. We say that $M$ is a left $R$-module if there is a map
    \[
        R \times M \to M, \quad (r, m) \mapsto rm,
    \]
    such that
    \[
        \begin{aligned}
            r(m+n) & = rm + rn, \\
            (r+s)m & = rm + sm, \\
            (rs)m  & = r(sm),   \\
            1m     & = m
        \end{aligned}
    \]
    for all $r, s \in R$ and $m, n \in M$.

    A \textbf{right $R$-module} is defined similarly, using a map $(m, r) \mapsto mr$ and the axiom $m(rs) = (mr)s$. If $R$ is commutative, left $R$-modules and right $R$-modules coincide. In general, a right $R$-module may be viewed as a left $R^{\mathrm{op}}$-module.
\end{definition}

For each $r \in R$, the map $\varphi_r\colon M \to M$ given by $\varphi_r(m) = rm$ is a homomorphism of abelian groups. Let
\[
    \operatorname{End}(M) = \{\, \varphi\colon M \to M \mid \varphi \text{ is a homomorphism of abelian groups} \,\}.
\]
Then $\operatorname{End}(M)$ is a ring under composition.

\begin{theorem}
    Let $M$ be an abelian group. Then $M$ is an $R$-module if and only if there is a ring homomorphism
    \[
        \rho\colon R \to \operatorname{End}(M).
    \]
\end{theorem}

\begin{proof}
    $(\Rightarrow)$ Suppose $M$ is an $R$-module. Define $\rho\colon R \to \operatorname{End}(M)$ by
    \[
        \rho(r)(m) = rm \qquad (r \in R,\ m \in M).
    \]
    For each $r \in R$, the map $\rho(r)\colon M \to M$ is a homomorphism of abelian groups since
    \[
        \rho(r)(m_1+m_2) = r(m_1+m_2) = rm_1 + rm_2.
    \]
    Moreover, $\rho$ is a ring homomorphism:
    \begin{itemize}
        \item $\rho(r+s)(m) = (r+s)m = rm + sm = \rho(r)(m) + \rho(s)(m)$, so $\rho(r+s) = \rho(r) + \rho(s)$;
        \item $\rho(rs)(m) = (rs)m = r(sm) = \rho(r)(\rho(s)(m))$, so $\rho(rs) = \rho(r) \circ \rho(s)$;
        \item $\rho(1)(m) = 1m = m$, so $\rho(1) = \operatorname{id}_M$.
    \end{itemize}

    $(\Leftarrow)$ Suppose $\rho\colon R \to \operatorname{End}(M)$ is a ring homomorphism. Define scalar multiplication by
    \[
        (r, m) \mapsto \rho(r)(m).
    \]
    Then
    \[
        \begin{aligned}
            r(m_1+m_2) & = \rho(r)(m_1+m_2) = \rho(r)(m_1) + \rho(r)(m_2), \\
            (r+s)m     & = \rho(r+s)(m) = \rho(r)(m) + \rho(s)(m),         \\
            (rs)m      & = \rho(rs)(m) = \rho(r)(\rho(s)(m)) = r(sm),      \\
            1m         & = \rho(1)(m) = \operatorname{id}_M(m) = m.
        \end{aligned}
    \]
    Hence $M$ is an $R$-module.
\end{proof}

\begin{example}
    (1) A vector space over a field is a module over that field.

    (2) An abelian group is a $\mathbb{Z}$-module.

    (3) Any ideal of a ring $R$ is naturally an $R$-module.

    (4) Let $M$ be an abelian group. Then $M$ is an $\operatorname{End}(M)$-module via
    \[
        \operatorname{End}(M) \times M \to M, \quad (\varphi, m) \mapsto \varphi(m).
    \]

    (5) Let $V$ be a vector space over a field $K$, and let $A\colon V \to V$ be a linear transformation. Then $V$ is a $K[x]$-module via
    \[
        K[x] \times V \to V, \quad (f(x), v) \mapsto f(A)v.
    \]

    (6) If $|M| = 1$, then $M$ is called the zero module.

    (7) Let $R$ be a ring with $1$ and let $A \in M_n(R)$. Then the solution set of $AX=0$, where $X \in R^n$, is a right $R$-module.
\end{example}

\begin{definition}[$R$-submodule]
    Let $R$ be a ring with $1$, and let $M$ be a left $R$-module. A subset $N \subseteq M$ is called an \emph{$R$-submodule} of $M$ if
    \begin{enumerate}
        \item $0 \in N$;
        \item $n_1, n_2 \in N$ implies $n_1+n_2 \in N$;
        \item $n \in N$ implies $-n \in N$;
        \item $r \in R$ and $n \in N$ imply $rn \in N$.
    \end{enumerate}
    Equivalently, $N$ is a subgroup of the abelian group $(M,+)$ that is closed under the $R$-action.
\end{definition}

\begin{example}
    (1) If $V$ is a vector space over a field $K$, then every subspace of $V$ is a $K$-submodule of $V$.

    (2) If $G$ is an abelian group, then every subgroup of $G$ is a $\mathbb{Z}$-submodule of $G$.

    (3) The left regular module ${}_R R$ has precisely the left ideals of $R$ as its submodules.

    (4) Let $G$ be an abelian group, viewed as an $\operatorname{End}(G)$-module. Then the $\operatorname{End}(G)$-submodules of $G$ are exactly the fully invariant subgroups of $G$.
\end{example}

\begin{remark}
    A natural question is the following: are submodules of finitely generated modules again finitely generated? In general the answer is no. This property holds for all submodules if and only if the ring is Noetherian.
\end{remark}

\begin{example}[Submodule of a finitely generated module that is not finitely generated]
    Let $k$ be a field, and consider the polynomial ring
    \[
        R = k[x_1, x_2, x_3, \dots]
    \]
    in infinitely many variables.
    \begin{itemize}
        \item The ring $R$ is a cyclic left $R$-module, generated by $1$.
        \item Let $I = \langle x_1, x_2, x_3, \dots \rangle$. Then $I$ is a submodule of the left $R$-module $R$.
        \item The submodule $I$ is not finitely generated. Indeed, suppose
              \[
                  I = \langle f_1, \dots, f_n \rangle
              \]
              for some polynomials $f_1, \dots, f_n \in R$. Let $N$ be the largest index of any variable appearing in any of the $f_i$. Then every element of $\langle f_1, \dots, f_n \rangle$ lies in the subring $k[x_1, \dots, x_N]$, but $x_{N+1} \in I$ does not. This is a contradiction.
    \end{itemize}
    Thus $I$ is a submodule of the finitely generated $R$-module $R$ which is not finitely generated.
\end{example}

\begin{definition}
    Let $M$ be an $R$-module, and let $(M_i)_{i \in I}$ be a family of submodules of $M$. Then
    \[
        \bigcap_{i \in I} M_i
    \]
    is an $R$-submodule of $M$. Moreover,
    \[
        \sum_{i \in I} M_i
        =
        \left\{ \sum_{i \in I}^{\mathrm{finite}} m_i \,\middle|\, m_i \in M_i \right\}
    \]
    is also an $R$-submodule of $M$.
\end{definition}

\begin{definition}
    Let $M$ be an $R$-module, and let $S \subseteq M$. The submodule generated by $S$ is
    \[
        \langle S \rangle
        =
        \left\{ \sum_{j=1}^n a_j s_j \,\middle|\, n \ge 0,\ a_j \in R,\ s_j \in S \right\}.
    \]
    Equivalently,
    \[
        \langle S \rangle
        =
        \bigcap_{\substack{N \le M \\ S \subseteq N}} N.
    \]
    If $M = \langle S \rangle$, then we say that $M$ is generated by $S$. In particular, if $|S| = 1$, then $M$ is called \emph{cyclic}; if $|S| < \infty$, then $M$ is called \emph{finitely generated}.
\end{definition}

\begin{definition}
    Let $R$ be a ring, and let $M$ be a left $R$-module. We say that $M$ is \emph{Noetherian} if every ascending chain of submodules
    \[
        M_1 \subseteq M_2 \subseteq M_3 \subseteq \cdots
    \]
    stabilizes; that is, there exists $n \ge 1$ such that
    \[
        M_n = M_{n+1} = M_{n+2} = \cdots.
    \]
    Equivalently, $M$ is Noetherian if every submodule of $M$ is finitely generated.
\end{definition}



